\par\Needspace{6\baselineskip}\medskip
\SourceLeaf{149}\pdfbookmark[2]{107.}{article-107}
\begin{center}107.\end{center}
Denotando per $\chi$ arcum quemcunque inter $-180^0$ et $+180^0$ situm, et per $s$ sinum arcus $\tfrac12\chi$, constat esse
\[
\tfrac12\chi=s+\tfrac13\cdot\tfrac12s^3+\tfrac15\cdot\frac{1\mathpunct{.}3}{2\mathpunct{.}4}s^5+\tfrac17\cdot\frac{1\mathpunct{.}3\mathpunct{.}5}{2\mathpunct{.}4\mathpunct{.}6}s^7+\text{etc.}
\]
\noindent Porro fit
\[
\tfrac12\sin\chi=s\BareRoot{(1-ss)}=s-\tfrac12s^3-\frac{1\mathpunct{.}1}{2\mathpunct{.}4}s^5-\frac{1\mathpunct{.}1\mathpunct{.}3}{2\mathpunct{.}4\mathpunct{.}6}s^7-\text{etc.}
\]
\noindent adeoque
\[
\begin{aligned}
\chi-\sin\chi=4\Bigl(&\tfrac13s^3+\tfrac15\cdot\tfrac12s^5\\[.5em]
&+\tfrac17\cdot\frac{1\mathpunct{.}3}{2\mathpunct{.}4}s^7+\tfrac19\cdot\frac{1\mathpunct{.}3\mathpunct{.}5}{2\mathpunct{.}4\mathpunct{.}6}s^9+\text{etc.}\Bigr).
\end{aligned}
\]
\noindent Substituimus in hac serie pro $s$ deinceps $\tfrac12\BareRoot{\dfrac{r+r'-\rho}{a}}$ et $\tfrac12\BareRoot{\dfrac{r+r'+\rho}{a}}$, quaeque inde proveniunt multiplicamus per $a^{\frac32}$; ita respective oriuntur series
\[
\begin{aligned}
&\tfrac16(r+r'-\rho)^{\frac32}+\tfrac1{80}\cdot\frac1a(r+r'-\rho)^{\frac52}\\[.5em]
&\qquad+\tfrac3{1792}\cdot\frac1{aa}(r+r'-\rho)^{\frac72}+\tfrac5{18432}\cdot\frac1{a^3}(r+r'-\rho)^{\frac92}+\text{etc}
\end{aligned}
\]
\[
\begin{aligned}
&\tfrac16(r+r'+\rho)^{\frac32}+\tfrac1{80}\cdot\frac1a(r+r'+\rho)^{\frac52}\\[.5em]
&\qquad+\tfrac3{1792}\cdot\frac1{aa}(r+r'+\rho)^{\frac72}+\tfrac5{18432}\cdot\frac1{a^3}(r+r'+\rho)^{\frac92}+\text{etc.,}
\end{aligned}
\]
\noindent quarum summas denotabimus per $T$, $U$. Iam nullo negotio patet, quum sit $2\sin\tfrac12\delta=\pm\BareRoot{\dfrac{r+r'-\rho}{a}}$, signo superiori vel inferiori valente prout $2f$ infra vel supra $180^0$ est, fieri $a^{\frac32}(\delta-\sin\delta)=\pm T$, signo perinde determinato. Eodem modo si pro $\epsilon$ accipitur valor minor infra $180^0$ situs, fiet $a^{\frac32}(\epsilon-\sin\epsilon)=U$; accepto vero valore altero, qui est illius complementum ad $360^0$, manifesto fiet $a^{\frac32}(\epsilon-\sin\epsilon)=a^{\frac32}\mathpunct{.}360^0-U$. Hinc itaque colliguntur duo valores pro tempore $t$:
\[
\frac{U\mp T}{k}\qquad\text{atque}\qquad\frac{a^{\frac32}\mathpunct{.}360^0}{k}-\frac{U\pm T}{k}\mathpunct{\cdot}
\]
